\chapter{Energy Track: From Alpha to Forecasting}
\label{ch:energy}

The energy track set out to repeat the equity result on a physically
grounded graph: twenty European bidding zones connected by the actual
cross-border transmission network, with day-ahead spot prices from the
ENTSO-E Transparency Platform \cite{entsoe-tp}. \Cref{fig:interconnector}
shows that graph as configured in the repository --- twenty zones, 38
physical borders and HVDC links, reference data rather than anything
estimated from prices. The track ended somewhere more
interesting than its plan. The cross-sectional alpha framing failed --- not quietly,
but through three nested evaluation artifacts, each of which initially
produced spectacular numbers and each of which the project's own controls
caught (\cref{sec:en-negcontrol,sec:en-artifact,sec:en-redesign,sec:en-clip}).
The track was then reframed as a forecasting problem, where the graph
finally showed genuine, control-validated value, concentrated exactly
where relational structure should matter most: on cross-border spreads
(\cref{sec:en-forecast,sec:en-edge}). This chapter reports the arc in
order, because the order is the contribution. Two research questions run
through it: whether one kernel transfers to a physical graph without
manufacturing false positives (RQ2), and whether that graph improves
forecasts (RQ3).

\begin{figure}[t]
  \centering
  \includegraphics[width=0.78\textwidth]{interconnector.pdf}
  \caption{The energy graph: 20 European bidding zones and their 38
  physical cross-border interconnections (land borders and HVDC
  cables), drawn from the repository's reference data. Unlike the
  equity correlation graph, nothing here is estimated --- the topology
  is the grid. DE\_LU, with eleven interconnectors, is the network's
  natural hub.}
  \label{fig:interconnector}
\end{figure}

\section{One kernel, two graphs --- and a second negative control}
\label{sec:en-negcontrol}

The dual-track design goal was engineering economy: the equity GAT, its
training loop, its four gates, and its A/B anchors run unchanged on the
energy panel; only the graph (physical interconnectors instead of
estimated correlations), the label horizon (24 hours instead of 5 days),
and the node set (bidding zones instead of symbols) differ. That goal was
met literally --- the same kernel functions are called by both tracks ---
and the first energy run doubled as a control.

On synthetic power data (2{,}857 hourly snapshots, 57{,}140 rows,
generated without any cross-zonal coupling) the pipeline correctly finds
nothing: OOS IC around $-0.01$ in both retrain arms, and, crucially, the
attention value-add is \emph{negative} ($-0.74$ single, $-0.59$
walk-forward). When there is no relational signal, the GAT is worse than
uniform neighbour averaging, not better. This is the energy analogue of
the equity negative control, on a structurally different graph: even a
learned attention layer over a physical topology does not manufacture
signal from noise.

\section{Real data, implausible result}
\label{sec:en-artifact}

The live ENTSO-E run covered all twenty zones for the full year 2024 at
hourly resolution --- 175{,}200 rows, with every configured zone
returning data. Real power prices immediately justified the label's
engineering: they spanned $-427$ to $1{,}896$~EUR/MWh, negative in 3.95\%
of hours. The pipeline ran end to end without kernel changes. The
numbers it produced were an OOS IC of 0.20, an OOS Sharpe of 8.25, a
composite roughly one hundred times its best single input, and four
gates out of four.

The same kernel earns a Sharpe of about 1.4 on equities. A result this
much better is not a discovery; it is a symptom. Diagnosis found three
compounding artifacts. First, overlapping labels: a 24-hour forward
return sampled hourly gives consecutive observations that share 23 of
their 24 hours --- measured lag-1 autocorrelation 0.91 --- while the
Sharpe computation treats them as independent, understating the
denominator by an order of magnitude. Second, the ``signal'' is trivial:
a one-line predictor, the negative of the current spot price, scores a
cross-sectional IC of 0.235, \emph{higher} than the GAT's 0.20. Whatever
the model learned, it is dominated by deterministic diurnal
mean-reversion. Third, the annualisation constant was the equity one
(252 periods per year) applied to hourly data (8{,}760). On top of all
three sits a market-structure lookahead: the day-ahead auction publishes
all 24 hours of day $D$ at once on $D-1$, so an ``hourly'' 24-hour-ahead
label is partly contemporaneous information.

The point the report wants to preserve is that the magnitude itself was
the alarm. The same scepticism that made the synthetic controls a design
feature flagged the real-data number as not credible, rather than
publishing a Sharpe of 8.

\section{Redesign in the day-ahead frame}
\label{sec:en-redesign}

The redesign moved the evaluation to the market's actual decision
frequency. The day-ahead auction clears once per day, so the snapshot
became the delivery day: base-load-aggregated features from data up to
day $D$, a label equal to the genuinely-unknown next-day base-load price
change $D \to D{+}1$, non-overlapping by construction, annualised at 365.

The mechanical checks all came back clean. Label lag-1 autocorrelation
fell from 0.91 to $-0.10$; the shuffle-label control gave a validation IC
of 0.07 against 0.19 for real labels, so the remaining signal is not
leakage. And yet the result stayed implausible: GAT composite OOS Sharpe
11.08 --- with the trivial negative-price predictor at 6.15. The
interim reading was that the pipeline was now honestly measuring a real
but non-tradeable physical regularity: the law of one price across
interconnected zones, already priced by cross-border capacity auctions.
That reading was half right, and its other half is the most instructive
result in this report.

\section{The third artifact: a clip on the evaluation return}
\label{sec:en-clip}

A direct challenge --- \emph{have you actually confirmed the strategy
would lose nothing?} --- prompted one more check: recompute the same
strategies under different definitions of the evaluation return.
\Cref{tab:en-clip} shows what that found.

\begin{table}[htbp]
  \centering
  \caption{The return-definition diagnostic (real ENTSO-E daily frame,
  OOS). The $\pm0.8$ clip, inherited from the training label, was also
  applied to the evaluation return in E13 --- and manufactured the
  entire positive result.}
  \label{tab:en-clip}
  \begin{tabular}{lS[table-format=-2.2]S[table-format=-2.2]}
    \toprule
    Evaluation return & {trivial $-$price Sharpe} & {GAT Sharpe} \\
    \midrule
    floored + $\pm0.8$ clip (as trained) & +6.28 & +11.00 \\
    winsorised $\pm0.5$                  & +4.61 & {---}  \\
    log return                           & +3.75 & {---}  \\
    plain relative return (honest)       & -1.52 & -1.54  \\
    \bottomrule
  \end{tabular}
\end{table}

The mechanism is specific and worth stating plainly. The long-short
strategy shorts expensive zones. European power prices have fat right
tails --- scarcity spikes reaching $1{,}896$~EUR/MWh in this sample ---
and the $\pm0.8$ clip on the evaluation return caps exactly those
short-leg losses. Removing the clip turns a Sharpe-11 winner into a
strategy that loses money, for the GAT and the one-line baseline alike.
The rank-IC, meanwhile, stays at about 0.21 under both return
definitions: cheap zones really do tend to rise relative to expensive
ones, but ranking correctly is blind to tail magnitudes, and the tails
are where this strategy's money is lost. A positive IC coexisting with
negative PnL is the transferable lesson --- never transform the
evaluation return, and never read rank quality as profitability. The
equity track is unaffected: its evaluation return was never clipped.

The final verdict on the energy alpha question is therefore airtight and
negative. Under honest returns the day-ahead cross-sectional strategy
loses money (Sharpe about $-1.5$); every positive energy number along the
way was an artifact --- overlap, annualisation, then the return clip ---
and the GAT adds no tradeable value over a one-line baseline. The fix
(evaluation on unclipped returns) shipped as the default, so the
repository can no longer reproduce the inflated numbers.

\section{The reframe: does the graph improve forecasts?}
\label{sec:en-forecast}

What survives E13b is a genuine question the data can answer: not ``is
there alpha?'' but ``does the interconnector graph improve forecast
skill?'' --- measured with the ladder of \cref{sec:skill} on six months
of real ENTSO-E data (20 zones, 87{,}360 rows, with day-ahead load, wind
and solar forecasts, realised load, and generation mix as drivers).
\Cref{tab:energy-node-skill} shows the node-level ladder, and
\cref{fig:energy-ladder} draws it with one added ingredient per rung.

\begin{table}[htbp]
  \centering
  \small
  \caption{Node-level price forecast skill ladder (skill $=1-\mathrm{MSE}/\mathrm{MSE}_{\mathrm{persistence}}$, OOS, $k=24$h).}
  \label{tab:energy-node-skill}
  \begin{tabular}{lll>{\raggedright\arraybackslash}p{0.34\textwidth}}
    \toprule
    predictor & skill & rank ic & note \\
    \midrule
    persistence & 0.000 & 0.637 & reference (tomorrow = today) \\
    seasonal\_naive & 0.000 & 0.637 & diurnal carry \\
    no\_graph\_ridge & 0.224 & 0.636 & own physical drivers (ridge) \\
    uniform\_graph\_ridge & 0.355 & 0.584 & + interconnector-neighbour mean (unlearned graph) \\
    gat\_node & 0.347 & 0.612 & learned attention (PyG GATv2; 5-seed mean) \\
    gat\_congestion & 0.346 & 0.615 & + price-spread congestion edge feature (5-seed mean) \\
    \bottomrule
  \end{tabular}
\end{table}


\begin{figure}[htbp]
  \centering
  \includegraphics[width=0.54\textwidth]{energy_ladder.pdf}
  \caption{The node-level forecast ladder of
  \cref{tab:energy-node-skill}, drawn as a ladder. Each rung adds
  exactly one ingredient, so each skill increment isolates one cause.
  The published drivers enter at the ridge rung and are shared by every
  rung above it; the graph's contribution ($+0.131$) enters below any
  learning.}
  \label{fig:energy-ladder}
\end{figure}

Three findings, in decreasing order of robustness. First, the graph
carries real forecast value: uniform neighbour averaging over the
interconnector topology lifts skill by $+0.131$ over the no-graph ridge
with the same drivers. The synthetic negative control (independent
zones) shows no such lift, so the gain is genuine cross-zonal coupling.
This rung is pure numpy --- no learned model is involved --- which makes
it the hardest number in the chapter to argue with.

Second, learned attention beats the uniform anchor on cross-sectional
\emph{ranking} in five seeds of five (rank-IC 0.61 against 0.58) but is
a wash on MSE skill (better in two seeds of five). Uniform
neighbour-averaging turns out to be near MSE-optimal on this graph --- a
hard bar --- so the GAT's robust edge here is order, not magnitude,
echoing the equity finding that attention's value is a modest tilt on
broad pooling.

Third, the congestion hypothesis --- that attention should down-weight
saturated borders if told how congested they are --- is not supported.
Two operationalisations of the edge feature were tried: the current
cross-border price spread (a leak-safe proxy), and ground-truth physical
flows over day-ahead transfer capacity fetched for all 38 borders. The
proxy's apparent gain existed only in a hand-rolled dense GAT
implementation ($+0.031$ skill, five of five seeds) and vanished under
the standard PyG GATv2 ($-0.002$, two of five); the ground-truth version
beat node-only attention in one seed of five. The hypothesis is not
refuted --- capacity data covered only 36\% of borders, and the most
coupled region uses flow-based coupling that a flow-over-capacity ratio
cannot express --- but as a smooth additive edge feature, congestion
does not add forecast skill. The two-backend check deserves its
methodology note: a single-seed PyG run had initially shown $+0.072$, a
lucky draw that a five-seed rerun corrected, and the cross-implementation
comparison is what stopped a fragile effect from becoming a claim.

\section{Edge-level spread prediction: the relational payoff}
\label{sec:en-edge}

The last experiment moved the target instead of the model. A
cross-border price spread $\mathrm{spot}_a - \mathrm{spot}_b$ is
irreducibly relational --- it does not exist for a single node, and it is
the object that transmission rights actually price. A GAT computes node
embeddings by message passing over the interconnector graph; a small MLP
head predicts each border's next-day spread from the two endpoint
embeddings and the current spread. The baseline that matters is a ridge
that already sees both endpoints' drivers and the current spread ---
so any gain over it is whole-network context, not information the
baseline lacks. \Cref{tab:energy-edge-skill} shows the ladder over the
38 borders.

\begin{table}[htbp]
  \centering
  \small
  \caption{Edge-level cross-border spread prediction (38 borders, OOS, $k=24$h, 5-seed means).}
  \label{tab:energy-edge-skill}
  \begin{tabular}{lll>{\raggedright\arraybackslash}p{0.34\textwidth}}
    \toprule
    predictor & skill & rank ic & note \\
    \midrule
    edge\_persistence & 0.000 &  & spread today (reference) \\
    edge\_ridge & 0.192 & 0.67 & both endpoints' day-ahead drivers + current spread (no graph) \\
    edge\_gat & 0.248 & 0.68 & GAT node embeddings $\to$ edge MLP head (message passing; 5-seed mean) \\
    \bottomrule
  \end{tabular}
\end{table}


The edge GAT beats the both-endpoint ridge by $+0.056$ skill --- about
29\% relative --- in five seeds of five, and the synthetic control
inverts as it should: on independent zones the same architecture is
\emph{worse} than the ridge, so the real-data gain is structure, not
capacity. This is the project's strongest relational result, and its
placement is the point: node-level price forecasting barely needed the
graph, congestion-as-edge-feature was null, and the GNN's value
concentrates precisely on the target that cannot be defined without the
relation. Caveats stated once: a single six-month window, untuned
hyperparameters, and modest absolute skill --- spreads are noisy
differences.

\section{Summary of the energy evidence}
\label{sec:en-summary}

\begin{table}[htbp]
  \centering
  \small
  \caption{Energy forecasting track: what is and is not supported (real ENTSO-E data, 20 zones, 5-seed means).}
  \label{tab:energy-findings}
  \begin{tabular}{>{\raggedright\arraybackslash}p{0.24\textwidth}>{\raggedright\arraybackslash}p{0.22\textwidth}lll}
    \toprule
    finding & comparison & skill delta & seeds & verdict \\
    \midrule
    Graph helps price level & uniform - no-graph & +0.131 & control-validated & robust \\
    Attention vs uniform (ranking) & gat rank\_ic - uniform & +0.031 & 5/5 & robust \\
    Attention vs uniform (skill) & gat - uniform & -0.008 & 2/5 & wash \\
    Congestion edge feature (price-spread proxy) & gat\_congestion - gat\_node & -0.002 & 2/5 & null \\
    Congestion edge feature (real flow/NTC) & flow - node & -0.005 & 1/5 & null \\
    Edge-level spread (message passing) & edge\_gat - edge\_ridge & +0.056 & 5/5 & robust * \\
    \bottomrule
  \end{tabular}
\end{table}


\Cref{tab:energy-findings} condenses the track. As an alpha project the
energy track concluded negatively three artifacts deep, and the honest
number --- the strategy loses money --- is stated as the finding it is.
As a forecasting project it produced two robust, control-validated
positives (the graph lift of $+0.131$ at node level; the edge-level
message-passing gain of $+0.056$, five of five seeds) and one clean null
(congestion edge features, under both operationalisations). As
methodology it produced the report's most transferable lessons: evaluate
on untransformed returns, distrust magnitudes that outrun their
plausibility class, and require cross-implementation replication before
believing a learned effect.
